% Immutable cited-input supplement. Target proofs remain embargoed.
% Target: night:P4.31; tag 02X4.
% Frozen source SHA256 74497bfe327fae25c9cdcb17334b3bf42a9793a9df7c0a4404d97568f2f760f0.
% CompendiumRef is resolved against the cited S6_09 source edition, not silently identified with the current series statement.

% Exact cited input: series:II / night:P4.22
\begin{proposition}[P4.22]\hypertarget{tag-02WU}{}\label{night:P4.22}\label{stmt:P4.22}


In the standing torus-fibration setting of this section, an irreducible singular support
$S$ whose smooth normalization $T$ is birationally ruled
over an elliptic curve has $e(S)\ge0$. Equality holds
exactly when $T$ is minimal.
\end{proposition}

% Exact cited input: series:II / night:P4.29
\begin{proposition}[P4.29]\hypertarget{tag-02X1}{}\label{night:P4.29}\label{stmt:P4.29}


Let $S$ be an irreducible singular support
in Part~4.2 whose normalization has ruled
resolution, and suppose $e(S)=0$.
In the double-horizontal-divisor case,
when the generic transverse type is
cuspidal $A_2$, the normalization is
a smooth minimal elliptic ruled surface.
The image conductor is a smooth elliptic
curve and every point along it has the
product cusp germ of P4.28.
\end{proposition}

% Exact cited input: series:II / night:P4.30
\begin{proposition}[P4.30]\hypertarget{tag-02X3}{}\label{night:P4.30}\label{stmt:P4.30}


Let $C$ be a compact irreducible reduced curve in the smooth
threefold $Y$ of Part~4.2. Suppose its normalization is an elliptic
curve $E$, and every point of $C$ is unibranch. Write $P$ for its
ideal in $Y$ and $P^{(2)}$ for the ideal of functions vanishing to
order at least two at the generic point of $C$. Then
\[
 \chi(\mathcal O_Y/P^{(2)})\le0.
\]
Equality forces $C$ to be smooth. The assertion uses
$c_1(TY)\cdot C=0$, which follows from $b_2(Y)=0$.
\end{proposition}

% Exact cited input: series:II / night:P4R.6
\begin{proposition}[P4R.6]\hypertarget{tag-02YI}{}\label{night:P4R.6}\label{stmt:P4R.6}


Let $S$ be an irreducible singular support in
Part~4.2, let $\nu:N\to S$ be its normalization,
and let $\rho:T\to N\to S$ use the minimal
resolution of $N$. There is an effective integral
divisor $D$ on $T$ such that
\[
 \omega_T(D)\simeq\rho^*\omega_S,
 \qquad \operatorname{Supp}D
       =\rho^{-1}(\operatorname{Sing}S)_{\rm red}.
\]
In particular $K_T+D$ is torsion, by P4.12.
The strict-transform conductor has its usual
positive integral coefficients. Every connected
exceptional configuration over a singular point
of $N$ has strictly positive coefficients in $D$.
Every such configuration meets the strict-transform
conductor.
\end{proposition}

% Exact cited input: series:II / night:P4R.7
\begin{proposition}[P4R.7]\hypertarget{tag-02YJ}{}\label{night:P4R.7}\label{stmt:P4R.7}


In Part~4.2, every irreducible singular
support whose normalization has a
resolution birationally ruled over a
curve of genus at least two has
nonnegative Euler number. The
normalization may be singular.
\end{proposition}

% Exact cited input: series:II / night:P4R.9
\begin{proposition}[P4R.9]\hypertarget{tag-02YL}{}\label{night:P4R.9}\label{stmt:P4R.9}


In Part~4.2, every irreducible singular
support whose normalization has a
resolution birationally ruled over an
elliptic curve has nonnegative Euler
number. When the horizontal part of
the pole divisor on the minimal ruled
model is reduced, the normalization
is smooth and minimal and the support
has Euler number zero.
\end{proposition}

